\documentclass[conference]{IEEEtran}
\usepackage{booktabs}
\usepackage[hidelinks]{hyperref}

\begin{document}

\title{D1. GraphRAG vs Vector RAG, Refereed Properly\\Midway Check --- 12 November 2026}

\author{\IEEEauthorblockN{Safina Mukhamedshina}
\IEEEauthorblockA{s.mukhamedshina@innopolis.university}
\and
\IEEEauthorblockN{Karina Shaikhutdinova}
\IEEEauthorblockA{k.shaikhutdinova@innopolis.university}
\and
\IEEEauthorblockN{Artem Mikhailin}
\IEEEauthorblockA{a.mikhailin@innopolis.university}
}

\maketitle

\noindent\textbf{Repository:} \url{https://github.com/safiction/graph-vs-vector}

\section{Implemented work}
The four-system comparison uses one frozen EnterpriseRAG-Bench subset:
15,000 documents and 45,142 shared chunks. BM25, BGE vector retrieval,
GraphRAG-local/global and the shared answer generator are implemented.
Full graph extraction and community summaries are complete, with a passing
saved structural provenance audit; this does not certify semantic correctness.
Corpus manifests, model settings and run logs support reproducibility;
large data and indexes are distributed separately from Git.

The baselines produced 24 development answers (12 per system) and 80
preliminary answers (40 per system, ten questions per type). All calls
succeeded, but eight preliminary answers reached the 384-token output cap.
The generator is Qwen2.5 3B Q4\_K\_M with a shared 2,560-token context cap.
Graph retrieval traces exist, but graph-generated answers are still missing.

\section{Preliminary results}
Table~\ref{tab:results} separates final-context document coverage on the
40-question preliminary evaluation set from manually graded correctness on
the separate development set. Correctness is the per-answer fraction of
satisfied reference facts, averaged across answers.

\begin{table}[ht]
\centering
\caption{Baseline results (\%). Coverage and correctness use different question sets.}
\label{tab:results}
\small
\setlength{\tabcolsep}{4pt}
\begin{tabular}{@{}lrrrr@{}}
\toprule
& \multicolumn{2}{c}{Prelim. coverage} & \multicolumn{2}{c}{Dev. correctness} \\
Type & BM25 & Vector & BM25 & Vector \\
\midrule
Single-fact   & 100.00 & 90.00 & 100.00 & 50.00 \\
Multi-hop     & 80.00 & 67.50 & 65.00 & 65.83 \\
Corpus-global & 68.75 & 52.08 & --- & --- \\
Fuzzy         & 80.00 & 70.00 & 50.00 & 0.00 \\
\bottomrule
\end{tabular}

\smallskip
\parbox{\columnwidth}{\footnotesize
Coverage: ten questions per type, except corpus-global, where two lack
gold document IDs and are N/A (eight applicable). Development correctness:
four questions per available type per system; no corpus-global examples.}
\end{table}

Safina and Karina independently graded all 24 development answers,
68 facts each. Initial fact-level agreement was 94.12\%, with Cohen's
$\kappa=0.880$. Four disagreements were adjudicated; none remain unresolved.
This is human--human agreement, not judge calibration. These small development
results do not establish a final winner or a learned routing policy.
Values trace to the baseline coverage exports and
\path{reports/preliminary/development_grading.json}.

\section{Blockers and immediate next steps}
\begin{enumerate}
\item \textbf{Graph answers (Artem/Karina):} run the completed graph through
the shared generator with identical settings and question IDs. Missing:
24 development and 80 preliminary answers. Only 80 of the planned 160
four-system preliminary attempts are currently available.
\item \textbf{Evaluation (Safina/team):} grade the remaining outputs, resolve
corpus-global development coverage, and agree aggregate failure/N/A handling.
No automated judge has been calibrated; balanced 40-answer, four-type
calibration and the paired four-system results table remain incomplete.
\item \textbf{Routing/compute (Artem/team):} fit routing only after paired
development scores are available; demonstrate T4 feasibility and consolidate
indexing costs, including recovery attempts. Full core runs, uncertainty
intervals and final analysis remain outstanding.
\end{enumerate}

\section{Changes from the proposal}
Following course feedback, we replaced full-corpus extraction with the fixed
15k subset. Selection is source-proportional and includes question gold
documents; this conditioning and possible changes to corpus-global answer
completeness must be disclosed. Raw-source retrieval is reported separately from generated
GraphRAG-global summaries.

The target remains 200 evaluation questions, 50 per type. Project-related
questions now belong only to multi-hop; the corpus-global allocation was
adjusted accordingly. The revised protocol selects routing on development
data before held-out evaluation, replacing the proposed two-fold routing
cross-validation. The current 40 preliminary questions remain evaluation
questions and must not be reused for tuning.

\end{document}
